% do not change this
\pdfbookmark[0]{\IfLanguageName{ngerman}{Zusammenfassung}{Abstract}}{abstract}
\begin{center} 
    \huge \IfLanguageName{ngerman}{Zusammenfassung}{Abstract}
\end{center}

% change from here

Die Seminararbeit untersucht die Log4Shell-Sicherheitslücke in der Open-Source-Bibliothek Apache Log4j. Diese Schwachstelle, auch bekannt als CVE-2021-44228, erlaubt es Angreifern, beliebigen Code auf betroffenen Systemen auszuführen. Dies geschieht durch speziell gestaltete Anfragen. Die Arbeit beginnt mit einer Einführung in die Funktionen und Hauptkomponenten von Log4j. Danach folgt eine detaillierte Analyse der Entdeckung und Ausnutzung der Schwachstelle. Es wird gezeigt, wie Angreifer JNDI-Injektionen nutzen, um die Lücke auszubeuten. Ein Beispielangriff wird ebenfalls erläutert. Die Arbeit geht auch auf die Auswirkungen auf Betriebsprozesse und Ressourcen ein. Zudem werden die langfristigen Bedrohungen und Herausforderungen durch Log4Shell diskutiert. Abschließend werden Strategien zur Prävention und Behebung vorgestellt. Diese sollen Unternehmen helfen, sich in Zukunft gegen ähnliche Bedrohungen zu schützen.